\section{Local mass and equations}\label{loc:section}

We retain the fixed data of \Cref{geo:setup,inc:markings,ht:conventions}.
The center is denoted by $t_0$, its labeled pair by
$\mathbf B=(B_1,B_2)$, and the chosen pair at a selected target $t$ by
$\mathbf P=(P_1,P_2)$.  Throughout this section, $h=h(t)\geq2$ and
$d=\max\{2,[K(t):K]\}$.  We write $C$ for positive constants depending
only on the fixed data, allowing their value and exponent to increase.
In particular, $Cd^C$ always denotes a polynomial bound with fixed
constants.  The endomorphism pair is supplied by
\Cref{ht:endomorphisms}, and its coordinate estimates, including the
versions with fixed additional denominators, are those of
\Cref{ht:coordinates}.

The argument has three steps.  The divisor $[t_0]$ supplies a large
weighted sum of local proximity to its own branch.  At the retained
finite places, good-reduction comparison converts horizontal transport
into traces of actual endomorphisms of a common reduction.  Grouping
these equations then supplies all the numerical hypotheses of
\Cref{int:height}.  The nonidentity hypothesis of that theorem is proved
separately in \Cref{nid:nonidentity}.

\subsection{The actual center and the amount of proximity}

Choose a rational parameter $x$ with a simple zero at $t_0$, regular on
the chosen affine neighborhood, and put $a=x(t)$.  The finite omissions
in \Cref{geo:setup,ht:conventions} ensure $a\ne0,\infty$.  Although $x$
may have other zeros on the completed curve, only its inverse branch
through $t_0$ will be used.  The same convention applies to every
conjugate of the fixed data.

All local sums below use the absolute-height weights of
\Cref{ht:conventions}.  For a field $L$ of data, these can be written
$w_v=[L_v:\Q_v]/[L:\Q]$ at places of $L$, or as the equivalent average
over embeddings into an algebraically closed local field.  Subsets of
embeddings are allowed.  Under a finite extension, the weights of the
embeddings lying above a given one sum to its original weight.  Thus
local extensions used to realize reductions, or a global normal closure
used to choose Frobenius, do not change any sum already defined.

Let $\mathcal S$ be a fixed finite set of rational primes containing the
primes required for the models, frames, levels and fixed centered data.
In particular, it contains the two auxiliary primes and the primes
dividing the multipliers of \Cref{geo:hecke-list}.  It will also contain
the finitely many primes specified in \Cref{loc:center-isogenies} below.
Outside $\mathcal S$, the center, its affine chart, parameter and
regular frame spread out smoothly, with an \'etale parameter and an
abelian scheme over that chart.  The inverse branch of each curve
coordinate at $t_0$ consequently has integral power-series
coefficients.  Write
\[
 \beta_p=\frac{\log p}{p-1}.
\]
At such a prime we retain only embeddings for which $t$ lies on that
actual center branch and
\begin{equation}\label{loc:radius}
 s_v:=-\log|a|_v>2\beta_p.
\end{equation}
At the archimedean places and at primes in $\mathcal S$, choose fixed
sufficiently small disks on the actual branches.  At a place where the
center is potentially good, these disks are chosen to satisfy the
good-reduction and convergence conditions below.  At a place where it
is not potentially good, \Cref{inc:bad-disks} excludes selected targets
from a sufficiently small disk, so no such place is retained.  Denote
the resulting collection of nearby embeddings by $V_0$.

\begin{lemma}[Proximity mass]\label{loc:mass}
For fixed constants $c>0$ and $C$, every selected target satisfies
\begin{equation}\label{loc:mass-estimate}
 \sum_{v\in V_0}w_v\log\frac1{|a|_v}
 \geq c h-C\bigl(1+\log(dh)\bigr)^2.
\end{equation}
The constants are uniform over arbitrary ramification of the target
fields.  They may depend on the fixed disks and on $\mathcal S$.
\end{lemma}

\begin{proof}
Let $D_0=[t_0]$ on the smooth projective completion of the selected
curve.  Its degree is positive.  Choose local Weil functions
$\lambda_{D_0,v}$ compatible with the fixed integral models away from
$\mathcal S$.  Their weighted sum is a height $h_{D_0}(t)$ up to a
bounded constant.  With the height convention using $D_0$, this is
$h+O(1)$; for any other fixed ample convention there is still a
constant $c>0$ with $h_{D_0}(t)\geq ch-O(1)$.  For the latter assertion,
choose a positive rational multiple of $D_0$ whose difference with the
other ample divisor has positive degree; the height of that difference
is bounded below.  This justifies the comparison without introducing an
error that grows linearly with $h$ on the wrong side.

On the disk through $t_0$, the ideal of $D_0$ is generated by $x$.
Hence
\[
 \lambda_{D_0,v}(t)=\log\frac1{|x(t)|_v}+O_v(1).
\]
At the spread-out primes the error is zero in the center tube, and the
local Weil function is zero outside that tube.  These assertions can
also be read from projective distance to the fixed point: on the
smooth chart its ideal is generated by the \'etale coordinate, whereas
outside its residue tube one of the ideal generators is a unit.
This uses the ideal of $t_0$, and never assigns proximity to another
zero of $x$ to $D_0$.

At each of the finitely many remaining local embeddings of the fixed
data, $\lambda_{D_0,v}$ is bounded above outside any fixed small center
disk.  Within that disk it differs from $-\log|x|_v$ by a fixed bound.
These bounds hold on algebraic points over every finite local
extension: they follow from fixed analytic units and fixed strict
inequalities defining the disks.  At the bad center places the
exclusion in \Cref{inc:bad-disks} therefore bounds the entire possible
contribution by a constant.  Shrinking the other fixed disks also loses
only a constant in the weighted sum.  Before imposing
\eqref{loc:radius}, the retained actual-branch contributions thus have
sum at least $ch-O(1)$.

It remains to bound the cost of that cutoff.  Take a field of degree
$D\leq Cd^C$ containing $a$ and the labeled target data, before taking
any normal closure.  We have $h(a)\leq Ch$ by
\Cref{ht:conventions}.  Let $\mathcal T(a)$ be the rational primes
admitting an embedding with $|a|_v<1$.  If $v\mid p$ is one such place,
then, with integral valuation $\ord_v$ and residue degree $f_v$,
\[
 w_v\log\frac1{|a|_v}
 =\frac{f_v\ord_v(a)}{D}\log p\geq\frac{\log p}{D}.
\]
Summing positive valuations and using the product formula gives
\begin{equation}\label{loc:prime-support}
 \sum_{p\in\mathcal T(a)}\log p\leq D h(a),
 \qquad N:=|\mathcal T(a)|\leq CDh.
\end{equation}
This argument sums the positive valuations themselves; cancellations
in the rational norm of $a$ are irrelevant.

The total weight over each rational prime is at most one.  The discarded
mass there is at most $2\beta_p$.  Order the $N$ primes increasingly as
$p_1,\ldots,p_N$.  The function $\log u/(u-1)$ decreases for $u>1$, and
$p_i\geq i+1$.  Consequently
\begin{equation}\label{loc:prime-loss}
 \sum_{p\in\mathcal T(a)}\beta_p
 \leq\sum_{i=1}^N\frac{\log(i+1)}i
 \leq C(1+\log(N+1))^2
 \leq C(1+\log(dh))^2.
\end{equation}
This proves \eqref{loc:mass-estimate}.  Both the support bound and the
weight identity survive field extension, so no ramification degree or
degree of a normal closure enters the estimate.
\end{proof}

\subsection{Good-reduction comparison on a common tube}

In the rational frame of $\cP_{\mathrm{dR}}$ let $Y(a)$ be horizontal
back-transport from $t$ to $t_0$, normalized by $Y(0)=1$.  Thus $Y(a)P$
is expressed in the fixed center frame.  This convention is used for
the primed data as well.

\begin{proposition}[Rational comparison and transport]\label{loc:comparison}
At every retained finite-place embedding there is, after a finite local
extension, a smooth proper abelian family over a smooth integral
moduli chart such that $A_{t_0}$ and $A_t$ have the same polarized
special fiber $\overline A$.  Rational comparisons of their homologies
with the cohomology of $\overline A$ identify their change of comparison
on $\cP$ with $Y(a)$.  In particular,
\begin{equation}\label{loc:common-traces}
 \langle B_i,Y(a)P_j\rangle
   =\Tr\bigl(\overline B_i\,\overline P_j\bigr)
 \qquad(1\leq i,j\leq2).
\end{equation}
Here the bars are actual specializations of the endomorphisms of the
two lifted fibers.  The same comparison is functorial for polarized
isogenies with specified reductions.  The assertion allows arbitrary
ramification of the local field of $t$.
\end{proposition}

\begin{proof}
Outside $\mathcal S$ the abelian family and the common specialization
are part of the integral chart already chosen.  At a retained prime in
$\mathcal S$, the center is potentially good.  Choose an auxiliary
fine level invertible in that residue characteristic and a finite
local extension over which the center has good reduction.  The center
then lies in the interior of the smooth integral moduli space.  A
sufficiently small fixed analytic neighborhood maps into the tube of
its special point.  Every algebraic point in that neighborhood, even
over a further ramified extension, therefore gives a good-reduction
lift with that same special fiber.  The disks defining $V_0$ have been
chosen inside these neighborhoods.

For clarity, the comparison used here is a comparison of rational
cohomology spaces.  For a complete mixed-characteristic DVR with
perfect residue field $k$, fraction field $L$, and a smooth proper
lift $\mathscr X$ of $X_0$, it has the form
\begin{equation}\label{loc:rational-comparison}
 H^1_{\mathrm{crys}}(X_0/W(k))\otimes_{W(k)}L
   \simeq H^1_{\mathrm{dR}}(\mathscr X_L/L).
\end{equation}
The rational comparison of Berthelot--Ogus
\cite{BerthelotOgus1983} is functorial and allows every ramification
index; this form is stated in \cite[\S1.2]{Berthelot1982} and explicitly
distinguished from integral comparison in
\cite[Introduction, equations~(2)--(3)]{AbhinandanYoucis2025}.
It does not assert preservation of integral lattices.  We use its
relative realization to compare the two lifts, as follows.

Embed the common proper special fiber $\overline A$ as the closed fiber
over the base special point in the smooth formal total family.  Its
rigid cohomology is computed by the de Rham complex on its tube;
properness eliminates a compactification boundary for this special
fiber.  Independence of the smooth embedding, functoriality and
extension of coefficients identify restriction to each smooth proper
lift with \eqref{loc:rational-comparison}; these are the tube
constructions of \cite[\S\S1.4--1.5 and 2.5(c)]{Berthelot1982}.
One may work on a smaller base disk, on which the relative
Gauss--Manin connection is ordinary.  An absolute de Rham
hypercohomology class on the total tube maps to a horizontal relative
class: the connecting morphism obtained by filtering the absolute
de Rham complex by base differential forms is precisely the Gauss--Manin
connection, and its composition with the absolute-to-relative map is zero.
The evaluation isomorphism with the cohomology of a proper lifted
fiber then shows that a basis of the common rigid cohomology gives a
basis of horizontal relative classes on this disk.  The change
between two evaluations is therefore the Taylor isomorphism; the
convergent Taylor formalism is described in
\cite[\S4.1]{Berthelot1982}.
All these operations commute with extension of the coefficient field,
so the same diagram holds for ramified evaluation fields.  The
unramified version of this diagram is displayed in
\cite[\S3.3, diagram~(3.9)]{LawrenceVenkatesh2020}.

After pullback to the parameter disk, the change of comparisons is
therefore horizontal and has identity value at $a=0$.  Uniqueness for
the ordinary differential equation identifies it with the Taylor
solution throughout the retained convergence disk.  Dualizing gives
the statement for covariant homology, and taking endomorphisms gives
the stated action on $\cP$.  Functoriality identifies an endomorphism
of either lift with its specialization, and likewise identifies an
isogeny with its specialized map.  The polarization and trace pairing
are preserved in these identifications.  Formula
\eqref{loc:common-traces} follows.

This argument uses a common good-reduction family.  The disks at a
center which is not potentially good have already been excluded by
\Cref{inc:bad-disks}.  Also, horizontal transport here identifies
cohomology and its tensors; no assertion about horizontal constancy of
the Hodge filtration is needed.
\end{proof}

\subsection{The three equations}

For two labeled pairs in a quadratic space, use the matrix notation
\[
 \langle\mathbf D,T\mathbf E\rangle
 :=\bigl(\langle D_i,TE_j\rangle\bigr)_{1\leq i,j\leq2}.
\]
Thus each matrix equation below is a system of four scalar equations.

At a good-reduction place, \Cref{loc:comparison} expresses this matrix
as rational traces. Using those trace values as coefficients at every
varying prime could give a number of distinct systems depending on the
target height. We instead compare the matrix with its simultaneous Frobenius
conjugate: conjugating both factors preserves each trace. For target
pairs on the fixed Hecke list, we will rewrite that equality using only
one horizontal transport. The next lemma makes this possible. A
fixed-multiplier isogeny between the targets specializes to an isogeny
between the center reductions; outside finitely many fixed primes, that
specialized map lifts to one of finitely many center isogenies.

\begin{lemma}[Bounded center isogenies under specialization]
\label{loc:center-isogenies}
Fix two of the conjugate centers and one multiplier $m_*$ from
\Cref{geo:hecke-list}.  There are finitely many characteristic-zero
polarized isogenies of multiplier $m_*$ between these centers.
After enlarging $\mathcal S$ by a fixed finite set, every polarized
isogeny of that multiplier between their good special fibers is the
specialization of one of these center isogenies.  This assertion is
simultaneous for the fixed finite list of centers and multipliers.
\end{lemma}

\begin{proof}
A polarized isogeny $\phi$ of multiplier $m_*$ satisfies
$\phi^*\lambda'=m_*\lambda$.  Since the fibers have dimension two,
$\deg\phi=m_*^2$, and its kernel lies in $A[m_*]$.  Away from primes
dividing $m_*$, kernels are among a fixed finite set of subgroup schemes
of this finite \'etale torsion scheme.  For each such kernel, the
quotient is fixed.  Polarized identifications of that quotient with the
second center form either an empty set or a torsor under its finite
polarized automorphism group.  This proves finiteness in characteristic
zero.

Here is a finite-type construction that also includes all special-fiber
possibilities.  After a fixed extension and localization, choose
relatively ample symmetric line bundles $L_0,L_0'$ inducing
$N\lambda_0,N\lambda_0'$ for the same fixed integer $N>0$.  On the
graph of any isogeny of multiplier $m_*$, the product bundle restricts
to $L_0\otimes\phi^*L_0'$, numerically equivalent to
$L_0^{\otimes(1+m_*)}$.  Riemann--Roch on the abelian surface gives
the fixed Hilbert polynomial
\[
 \chi\bigl(A_{t_0},(L_0\otimes\phi^*L_0')^{\otimes n}\bigr)
   =N^2(1+m_*)^2 n^2.
\]
All graphs therefore lie in one finite-type relative Hilbert scheme
of the product of the two fixed abelian schemes.  The locus on which
the first projection of the graph is an isomorphism is open.  On this
graph locus, preservation of the origin is closed and makes the
resulting morphism a homomorphism by rigidity for abelian schemes.
The condition $\phi^\vee\lambda_0'\phi=m_*\lambda_0$ is a closed
equality of homomorphisms to the dual abelian scheme.  This constructs
a separated finite-type scheme of the desired polarized homomorphisms;
the polarization identity forces each of them to be an isogeny.
The construction uses the same Hilbert polynomial in every residue
characteristic, so it includes any additional special-fiber polarized
automorphisms as well.

After a fixed extension making its finitely many geometric generic
points rational, finite presentation spreads each of them to a section
over an open arithmetic base.  A common localization works for all
of them.  Each section has closed image by separatedness.  The
complement of the finite union of these images is consequently open,
hence constructible, and has empty generic fiber.
Its image in the one-dimensional arithmetic base is constructible
and misses the generic point, hence consists of finitely many closed
points.  Removing their residue characteristics proves the assertion.
One can obtain the same conclusion by spreading out the finite kernel
list and the polarized-identification schemes.  Both constructions
involve fixed centers and fixed multipliers, so the removed primes are
independent of $t$ and of any target isogeny subsequently chosen.
\end{proof}

\begin{proposition}[Local equations]\label{loc:equations}
The nearby embeddings $V_0$ can be partitioned into groups using one of
the following three systems, together with the base-coordinate
equations in \eqref{loc:base-equations} below:
\begin{align}
 \langle\mathbf B,Y(a)\mathbf P\rangle
   &=\mathbf r,
   &&\mathbf r\in\Mat_{2\times2}(\Q),
   \label{loc:single}\tag{1}\\
 \langle\mathbf B,Y(a)\mathbf P\rangle
   &=\langle\mathbf B',Y'(aR)\mathbf P'\rangle,
   &&R=a'/a,
   \label{loc:paired}\tag{2}\\
 \langle\mathbf B,Y(a)\mathbf P\rangle
   &=\langle I_0^{-1}\mathbf B',Y(a)\widehat{\mathbf P}\rangle,
   &&\widehat{\mathbf P}=I_t^{-1}\mathbf P'.
   \label{loc:exception}\tag{3}
\end{align}
Here primes denote a simultaneous conjugate of the labeled data.
Type~\eqref{loc:single} is used at infinity and at retained primes in
$\mathcal S$.  At the other retained primes, type~\eqref{loc:paired}
is used if the coarse target pair lies outside the closed finite
Hecke list of \Cref{geo:hecke-list}, and type~\eqref{loc:exception}
is used if it lies on that list.  In types~\eqref{loc:paired} and
\eqref{loc:exception}, the following additional data and properties
hold at every assigned embedding:
\begin{enumerate}[label=(\roman*)]
\item $|a'|_v=|a|_v$, hence $|R|_v=1$, and $h(R)\leq Ch$.
\item There are two actual common reductions $\overline A$ and
$\overline A'$ of the unprimed and primed center--target pairs,
respectively.  Relative Frobenius
$F:\overline A\longrightarrow\overline A'$ has polarized multiplier
$p$ and conjugates both labeled pairs to their primed counterparts.
\item On each common reduction, smooth proper specialization at
$\ell_1\ne p$ preserves the two separately saturated plane lattices.
Their reductions span the independent degenerate four-space prescribed
by \Cref{inc:markings}.
\item In type~\eqref{loc:exception}, $I_0=\Ad(\phi_0)$ and
$I_t=\Ad(\phi_t)$ are conjugation isometries associated to polarized
isogenies of the same fixed multiplier $m_*$.  Their reductions satisfy
$\overline\phi_0=\overline\phi_t$, and
\begin{equation}\label{loc:transport-square}
 I_0Y(a)=Y'(a')I_t.
\end{equation}
Here $p\nmid\ell_1m_*$ and $\ell_1\nmid m_*$.  The maps induced by
$F$ and $\overline\phi_0$ on $\ell_1$-adic Tate modules are integral
isomorphisms, up to their unit similitude multipliers.
\end{enumerate}
The entries of $\mathbf r$ have bounded denominators and logarithmic
heights at most $C(1+\log h+\log d)$.  The pairs
$\widehat{\mathbf P}$ have fixed bounded denominators, polynomial
Rosati norms, and all the coordinate bounds of \Cref{ht:coordinates}.
\end{proposition}

\begin{proof}
At a retained finite place, \eqref{loc:common-traces} identifies the
left-hand entries with traces of products of actual specialized
endomorphisms.  For an integral endomorphism $f$ of an abelian variety,
the degree polynomial $\deg(n-f)$ is integral and equals the
characteristic polynomial on every prime-to-characteristic rational
Tate module; see
\cite[Proposition~12.9 and the preceding paragraph]{MilneAV}.
We explain the comparison with rational crystalline homology, since
the product $\overline B_i\overline P_j$ need not lift to either
characteristic-zero fiber. The cup-product map
$\bigwedge^r H^1_{\mathrm{crys}}(\overline A)\to
H^r_{\mathrm{crys}}(\overline A)$ is an isomorphism. Indeed,
the rational comparison with the center's good lift respects cup
product: the restrictions and base-change maps of the de Rham
complexes in the tube construction are multiplicative. It therefore
identifies the displayed map, after a faithful scalar extension, with the de Rham
cup-product map; that map is an isomorphism by comparison with the
cohomology of a complex torus. Being defined by cup product, the
crystalline isomorphism is natural for every endomorphism of
$\overline A$, including those that do not lift.

For an actual integral endomorphism $u$ of $\overline A$, the map
$[n]-u$ is an isogeny for every sufficiently large integer $n$.
Its pullback on first cohomology is $n-u^*$, by the K\"unneth
description of first cohomology and the group law. On top cohomology
its action is multiplication by $\deg([n]-u)$. One may verify this
last assertion in rigid cohomology, using the trace fundamental class
and compatibility of cycle classes with proper pushforward
\cite[\S2.1 and Proposition~6.4]{Petrequin2003}, and then use
smooth-proper crystalline--rigid comparison
\cite[\S1.5, Proposition~2]{Berthelot1982}. The degree statement
includes inseparable isogenies: it is the degree of a finite flat map,
not just the number of geometric kernel points; see also
\cite[Lemma~6.5]{Petrequin2003}.
Taking the top exterior power therefore gives
$\det(n-u^*\mid H^1_{\mathrm{crys}})=\deg([n]-u)$ for infinitely
many $n$, hence as a polynomial identity. Dualization preserves
the characteristic polynomial. Combining this with the Tate-module
identity proves that all the homological traces used here agree.
Thus these traces are rational and are
integers when both factors are integral; fixed rational denominators
introduce only fixed denominators in the trace.

All specialized operators here are symmetric for the common
specialized polarization.  Rosati positivity, valid in every
characteristic \cite[Theorem~17.3]{MilneAV}, and Cauchy--Schwarz give
\begin{equation}\label{loc:trace-size}
 \left|\Tr(\overline B_i\overline P_j)\right|^2
 \leq\Tr(\overline B_i^2)\Tr(\overline P_j^2)
 =\Tr(B_i^2)\Tr(P_j^2)\leq C(dh)^C.
\end{equation}
The absolute value in this inequality is the ordinary real absolute
value of the rational trace.  It is unrelated to the residue
characteristic norm.  At infinity, analytic Betti transport on the
fixed simply connected center disk identifies integral homology
lattices.  The analogous cross traces are again rational with fixed
denominators.  The Hodge metrics on this fixed disk, in a flat
trivialization, are uniformly comparable.  Cauchy--Schwarz for the
corresponding operator norms and the target Rosati bound give the
same polynomial size bound.  These observations establish
\eqref{loc:single} and its coefficient-height estimate at all of its
assigned places.

Now let $p\notin\mathcal S$.  Take a normal closure of a field
containing the labeled data, embed it at the assigned place, and choose
a lift of residue $p$-power Frobenius in its decomposition group.
The lift preserves this valuation.  Apply it simultaneously to the
center, parameter, target, frames, polarization, levels and both
endomorphism labels; denote the result by primes.  This gives
$|a'|_v=|a|_v$.  Absolute heights are invariant under conjugation, so
\[
 h(R)=h(a'/a)\leq h(a')+h(a)=2h(a)\leq Ch.
\]
The primed special fiber is the $p$-power twist of the unprimed one.
Functoriality of relative Frobenius gives, for the same isogeny $F$,
\[
 \overline B_i'=F\overline B_iF^{-1},\qquad
 \overline P_j'=F\overline P_jF^{-1}.
\]
These are conjugations in rational endomorphism algebras; no inverse
of Frobenius on an integral crystalline lattice is asserted.  Taking
the trace of the product removes $F$.  Applying
\Cref{loc:comparison} on each of the two actual common reductions
gives exactly \eqref{loc:paired}.

At these primes all fixed levels are invertible.  Integral smooth
proper specialization identifies the $\ell_1$-adic homology lattices
of each pair of lifts with the lattice of its common reduction,
compatibly with the specified level marking.  For either plane $H$
this transports the intersection of $H\otimes\Q_{\ell_1}$ with the
ambient integral operator lattice to the corresponding intersection
after specialization.  Thus the separately saturated lattices, and
their marked reductions, retain the incidence in
\Cref{inc:markings}.  This statement concerns the rational planes
and their intersections with lattices.  It does not require the
chosen small vectors $B_i,P_j$ themselves to be bases after reduction
modulo $\ell_1$.

If the coarse target pair belongs to the closed exceptional list,
choose a polarized isogeny $\phi_t:A_t\to A_{t'}$ of one of its fixed
multipliers $m_*$.  The list is fixed before the auxiliary levels;
only its pullback to those levels is being tested here.  After a local
extension, $\phi_t$ extends between the good-reduction abelian schemes.
Its reduction is an isogeny of multiplier $m_*$ between the same two
special fibers as the centers.  By \Cref{loc:center-isogenies}, it
equals the reduction of one of the fixed center isogenies $\phi_0$.
Conjugation by a polarized isogeny preserves trace and Rosati
symmetry, hence induces the isometries $I_t,I_0$.  Applying the
functorial comparisons to their identical reductions proves
\eqref{loc:transport-square}.  Substitution into
\eqref{loc:paired} gives \eqref{loc:exception}.

The identity $\phi_t^{-1}=m_*^{-1}\phi_t^\dagger$ shows that
$m_*\widehat P_j=m_*\phi_t^{-1}P_j'\phi_t$ is integral.  Since
conjugation preserves the Rosati form, the norms of
$\widehat P_j$ are bounded by the same polynomial as those of $P_j'$.
Apply \Cref{ht:coordinates} to these actual rational endomorphisms
of $A_t$ with that fixed denominator.  This supplies all their required
coordinate bounds directly; no height estimate for the entries of a
chosen isogeny matrix is needed.  Finally $p$, $m_*$ and $\ell_1$ are
pairwise coprime in the required positions, so the Tate-module
integrality and invertibility in (iv) follow from the isogeny degrees
$p^2,m_*^2$.
\end{proof}

The order of fixed choices is useful here.  The list of coarse Hecke
correspondences and its multipliers was chosen in
\Cref{geo:hecke-list}; the auxiliary levels were then chosen in
\Cref{inc:markings}.  The specialization lemma, or any later spreading
out or shrinking of charts, may enlarge $\mathcal S$.  A newly included
prime is handled by \eqref{loc:single} if the center is potentially
good there, and by \Cref{inc:bad-disks} otherwise.  Its contribution is
therefore either retained with a single trace equation or bounded
independently of the target.  No new multiplier, level or coarse
correspondence is chosen in this enlargement.

\subsection{Uniform grouping and the numerical interpolation data}

Fix affine coordinates $u_1,\ldots,u_b$ on the chosen curve chart, with
algebraic regular germs $f_\nu$ at the actual center so that
$u_\nu(t)=f_\nu(a)$ on its disk. Write $\mathbf u=(u_1,\ldots,u_b)$,
and use primes for the corresponding coordinates on a conjugate chart.
We now specify the coordinate tuple $U$ in each system:
\[
\begin{array}{c|c|c}
 \text{system} & U & \text{constant equation data} \\ \hline
 \eqref{loc:single}
   & (\mathbf u,\mathbf P) & \mathbf B,\mathbf r \\
 \eqref{loc:paired}
   & (\mathbf u,\mathbf u',\mathbf P,\mathbf P',R)
   & \mathbf B,\mathbf B' \\
 \eqref{loc:exception}
   & (\mathbf u,\mathbf P,\widehat{\mathbf P})
   & \mathbf B,I_0^{-1}\mathbf B'
\end{array}
\]
Every operator slot in $U$ contributes its coordinates in the fixed
rational five-dimensional frame. The listed constants remain fixed
within a group, as do its centered germs; $\mathbf r$ may depend on the
target. At the target the operator pairs have the endomorphism meanings
specified in \Cref{loc:equations}. On the ambient affine space their
slots are unrestricted coordinates. In particular, an algebraic variety
through the target used in \Cref{int:height} is not assumed to parametrize
Hodge endomorphisms. In the exceptional system the conjugation map $I_t$
enters only through the target value of $\widehat{\mathbf P}$.

Adjoin the scalar equations
\begin{equation}\label{loc:base-equations}
 u_\nu-f_\nu(a)=0\quad(1\leq\nu\leq b),
 \qquad
 u_\nu'-f_\nu'(aR)=0\quad\text{in type \eqref{loc:paired}}.
\end{equation}
These equations explicitly retain the curve projections and branches;
they will matter when testing identities on arbitrary algebraic
subvarieties in \Cref{nid:nonidentity}.

\begin{proposition}[Groups and numerical bounds]\label{loc:groups}
There are fixed constants $C$ such that either $h\leq Cd^C$, or there
is one system of type \eqref{loc:single}, \eqref{loc:paired} or
\eqref{loc:exception}, with \eqref{loc:base-equations}, having the
following properties.  Its target is an algebraic point $z=(a,U)$
in an affine space of fixed dimension, and it holds at all embeddings
in a weighted subcollection $V\subset V_0$.  Writing each scalar
expression as
\[
 E(a,U)=\sum_{j\geq0}a^j e_j(U),
\]
the numerical and function hypotheses of \Cref{int:height} hold:
\begin{align}
 h_{\mathrm{aff}}(z)&\leq Cd^C h,
 &\sum_{v\in V}w_v\log\frac1{|a|_v}&\geq\frac{h}{Cd^C},
 \label{loc:group-height-mass}\\
 \sum_{v\in V}w_v\log^+\|U\|_v
   &\leq Cd^C(1+\log h)^C,
 &\deg e_j&\leq C(1+j),
 \label{loc:group-local-degree}\\
 h_{\mathrm{aff}}(e_j)
   &\leq C(1+j+d)^C(1+\log h)^C,
 \label{loc:coefficient-height}\\
 \left|\sum_{j\geq l}a^j e_j(U)\right|_v
   &\leq |a|_v^l\exp\bigl((1+l)b_v\bigr),\notag\\
 \sum_{v\in V}w_v b_v&\leq Cd^C(1+\log h)^C.
 \label{loc:tail-bound}
\end{align}
Here $l\geq0$, $b_v\geq0$, and one common choice of $b_v$ works for
the bounded number of expressions.  Each expression is a polynomial
combination of fixed bounded degree of the coordinates, fixed regular
algebraic germs and fixed ordinary horizontal matrix germs, evaluated
at $a$ and, in type~\eqref{loc:paired}, at $aR$.  Only a fixed finite
list of such germs is used.  Constants in the expressions may depend
on the target, with the bounds displayed above.

There are at most $Cd^C$ groups in the construction.  This bound counts
ordered conjugate tuples, not elements of the Galois group of a normal
closure.  The proposition asserts no nonidentity condition on an
algebraic subvariety through $z$.
\end{proposition}

\begin{proof}
\emph{Counting the groups.}
Let $\tau$ comprise the fixed data, $t$ and the labeled endomorphism
coordinates, and choose a field of definition with absolute degree
$D\leq Cd^C$, as supplied by \Cref{ht:endomorphisms}.  There are at
most $D$ conjugate tuples.  At infinity and at the finitely many
primes in $\mathcal S$, group by the embedded tuple, the fixed branch
and its resulting rational matrix $\mathbf r$.  Each embedding gives
one such matrix, hence this uses at most $CD$ groups.  This counts
actual embeddings at fixed places; counting all rational matrices of
the allowed size would give an unnecessary dependence on $h$.

For the Frobenius equations, group by the ordered pair $(\tau,\tau')$.
There are at most $D^2$ values.  The normal closure in which Frobenius
was chosen can have much larger degree, but only these ordered tuple
values enter an equation.  The field generated by a pair has degree
at most $D^2$.  The primed fixed family belongs to the finite list of
conjugates of the fixed family; its germs are independent of the
varying target.

Split these groups according to whether the coarse pair is on the
closed exceptional list.  On it, refine by a choice of its fixed
multiplier, an isogeny $\phi_t$, and its specializing center isogeny
$\phi_0$.  The multiplier and center choices range over finite fixed
sets.  For a fixed multiplier, a target isogeny is specified by a
subgroup of bounded torsion and a polarized identification of the
quotient with the second target.  The number of subgroups is bounded
in dimension two and $m_*$.  The order of a polarized automorphism
group is uniformly bounded in fixed dimension: it is finite by
positivity, and acts faithfully on an integral homology lattice of
fixed rank, whose finite subgroups have bounded order.  Thus only a
fixed bounded relative extension is needed to choose $\phi_t$, and
the refinement still has at most $CD^2$ groups.  The transported tuple
also has polynomial degree.  Enlarging fields and lifting the
embedding weights in making this partition preserves their sum.

\emph{Selecting mass.}
By \Cref{loc:mass}, the total mass is at least
$ch-C(1+\log(dh))^2$.  Either $h\leq Cd^C$, or this is at least
$ch/2$, after increasing the fixed constants.  Indeed
$(\log h)^2\leq\varepsilon h+C_\varepsilon$ and
$(1+\log d)^2\leq Cd^2$ absorb the error.  In the latter case one of
the at most $Cd^C$ groups has mass at least $h/(Cd^C)$.  Fix that
group and call its embeddings $V$.  Its algebraic tuple and constants
give common equations at all of them.

\emph{Coordinate heights.}
Use the coordinate tuple $U$ specified above for the chosen system.
By \Cref{ht:coordinates}, the endomorphism coordinates have
global affine height $C(h+\log d)$ and weighted positive local logs
at most $C(1+\log h+\log d)$ at the chosen nearby embeddings.  This
also applies to $\widehat{\mathbf P}$ by
\Cref{loc:equations}.  Base-coordinate heights are $O(h)$, and their
positive local logs are bounded on the fixed center disks, with zero
cost away from the fixed primes.  We have already proved $h(R)\leq
Ch$, while $\log^+|R|_v=0$ on its group.  These facts give
\eqref{loc:group-height-mass} and the first bound in
\eqref{loc:group-local-degree}.

\emph{Fixed germs and their coefficients.}
The matrices $Y,Y'$ solve fixed ordinary algebraic differential
systems in the selected regular frames.  Dual or inverse matrices,
when needed to realize an operator system, satisfy fixed differential
systems as well.  Each $f_\nu$ is an algebraic regular germ on the
fixed chart.  There are only finitely many conjugate fixed charts and
centers.  Thus the number of systems, their ranks and the degrees of
polynomial combinations are fixed.  In particular, ``ordinary'' here
refers to the absence of a pole of the differential equation at its
center; it imposes no ordinarity condition on an abelian special fiber.

For completeness, consider a local system $Z'=A(x)Z$ at a spread-out
prime, with $A(x)=\sum_{i\geq0}A_i x^i$ integral and $Z(0)$ integral.
Writing $Z(x)=\sum_{n\geq0}Z_nx^n$ gives
\[
 (n+1)Z_{n+1}=\sum_{i=0}^n A_iZ_{n-i}.
\]
Induction proves $n!Z_n$ integral, since $n!/(n-i)!$ is an integer.
In the row convention for back-transport the same argument applies
with multiplication on the other side.  Therefore
\begin{equation}\label{loc:factorial}
 \|Z_n\|_v\leq |n!|_v^{-1}
 \leq \exp(n\beta_p).
\end{equation}
The last estimate follows from
$v_p(n!)=\sum_{r\geq1}\lfloor n/p^r\rfloor\leq n/(p-1)$.
Regular algebraic coordinate germs are integral in the \'etale chart.
A product of a fixed number of these series retains the factorial
bound: multiplying its coefficient of degree $n$ by $n!$ expresses it
as a sum with the integral factors $n!/(n_1!\cdots n_k!)$.

At each of the finitely many other places, ordinary-point convergence
gives fixed exponential coefficient bounds.  At the archimedean
embeddings this is Cauchy's estimate on a fixed smaller disk.  At a
fixed finite place it follows from the recurrence after rescaling the
parameter to make the analytic connection bounded.  Summing all local
coefficient bounds yields, for each fixed germ,
\[
 h_{\mathrm{aff}}(Z_n)\leq C(1+n\log(n+1)).
\]
Here the contribution from all factorial denominators is at most
$\log(n!)$, and the exceptional-place contribution is $O(n+1)$.
For the vector of coefficients up to index $j$ the same reasoning
gives a polynomial bound in $1+j$.  Fixed products and sums preserve
such a bound.

In $Y'(aR)$ the coefficient of $a^j$ is $Y_j'R^j$.  Thus the degree
in the polynomial variables grows at most linearly in $j$, without
altering its coefficient heights.  The same holds for $f_\nu'(aR)$.
The other coordinate factors occur with fixed bounded degree.  The
only varying constants of type~\eqref{loc:single} are $\mathbf r$,
whose coefficient heights are $C(1+\log h+\log d)$; the constants in
type~\eqref{loc:exception} are fixed centered isogeny data.  This
proves the remaining degree bound and
\eqref{loc:coefficient-height}.

\emph{Tails at the target.}
At $p\notin\mathcal S$, let $M_v\geq1$ bound all fixed-degree
coordinate and constant factors in the bounded collection of
expressions.  We may choose it so that
\[
 \log M_v\leq C\log^+\|U\|_v+C_v,
\]
where $C_v=0$ outside fixed primes; a larger fixed bound includes
constant terms when $l=0$.  The factorial argument gives
$|e_j(U)|_v\leq M_v\exp(j\beta_p)$.  The substitution $aR$ introduces
no further factor because $|R|_v=1$.  By
\eqref{loc:radius}, $|a|_v\exp(\beta_p)<1$, and the ultrametric
triangle inequality yields
\[
 \left|\sum_{j\geq l}a^j e_j(U)\right|_v
 \leq M_v\sup_{j\geq l}(|a|_v e^{\beta_p})^j
 =M_v|a|_v^l e^{l\beta_p}.
\]
Thus $b_v=\beta_p+\log M_v$ works, including $l=0$.
At the fixed places choose the evaluation disk with radius at most
one half of a fixed convergence radius in the archimedean case, and
strictly smaller in the nonarchimedean case.  The geometric-series
estimate gives the same bound with
$b_v=C_v+C\log^+\|U\|_v$ and, where necessary, an additional
$C\log^+\|\mathbf r\|_v$.  The sum of the latter contributions is
bounded by the coefficient height of $\mathbf r$ and the fixed number
of places.  The positive coordinate logs have already been bounded.
Finally \eqref{loc:prime-loss} bounds the sum of $\beta_p$ over every
prime in this group.  Enlarging $b_v$ to the maximum for the bounded
number of expressions proves \eqref{loc:tail-bound}.

The entire construction uses absolute weights.  Selecting a subgroup
of embeddings only decreases the nonnegative coordinate and tail
costs; field extension preserves the previous sums.  None of these
estimates requires a bound in terms of $d$ alone on the number of
supersingular, or of any other, proximity primes.  This completes the
numerical verification and leaves precisely the nonidentity condition
to \Cref{nid:nonidentity}.
\end{proof}
